%------------------------------------------------------------------------------%
\begin{question}
  Considerando a equação
  \[
    x^2-y^2-4x-2y+3=0
  \]
  \vspace{-\baselineskip}
  \begin{enumerate}
    \item identifique a curva
    \item determine seus principais valores e pontos
    \item esboce a curva
  \end{enumerate}
\end{question}

%------------------------------------------------------------------------------%
 \begin{resolution}{Completando quadrados}
   Agrupamos os termos

   \[
     x^2-4x-(y^2+2y)+3=0
   \]

   Completamos os quadrados

   \begin{align*}
     x^2-4x+4&=(x-2)^2
     \\
     y^2+2y+1&=(y+1)^2
   \end{align*}

   Substituindo na equação

   \begin{align*}
     (x-2)^2-4
     -(y+1)^2+1+3&=0
     \\
     (x-2)^2-(y+1)^2&=0
   \end{align*}
 \end{resolution}

 %------------------------------------------------------------------------------%
 \begin{resolution}{Classificação}
   Utilizamos a diferença de quadrados

   \begin{align*}
     (x-2)^2-(y+1)^2&=0
     \\
     [(x-2)-(y+1)][(x-2)+(y+1)]&=0
     \\
     (x-y-3)(x+y-1)&=0
   \end{align*}

   Portanto, a equação representa duas retas

   \[
     x-y-3=0
   \]

   \[
     x+y-1=0
   \]

   As retas se intersectam no ponto obtido pelo sistema

   \[
     \begin{cases}
       x-y=3\\
       x+y=1
     \end{cases}
   \]

   Somando as equações

   \[
     2x=4
     \qquad
     x=2
     \qquad
     y=-1
   \]

   Trata-se de uma cônica degenerada formada por duas retas concorrentes
 \end{resolution}

%------------------------------------------------------------------------------%
\begin{resolution}{Esboço}
  \centering
  \includegraphics{I-07-Identificacao-ex-F}
\end{resolution}

%------------------------------------------------------------------------------%
